% ===========================================================================
\section{Verificación}
\label{sec:verif}

El plan de verificación se escribió antes que el RTL (\repo{docs/spec/verificacion.md}). Fija
qué se prueba, contra qué se compara y cuándo una prueba pasa.

\subsection{Niveles de prueba}

\begin{center}
\small
\begin{tabular}{@{}l p{6.2cm} p{5.6cm}@{}}
  \toprule
  \textbf{Nivel} & \textbf{Qué prueba} & \textbf{Dónde} \\
  \midrule
  1. Unitario      & Un módulo aislado, con estímulos propios (barrido exhaustivo cuando la interfaz lo permite) & \repo{tb/unit/}, 27 testbenches \\
  2. Por etapa     & El pipeline hasta la etapa $N$, comparando los registros de segmentación con el modelo en cada ciclo & \repo{pipeline_tb.v} \\
  3. Core completo & Cada programa hasta el HALT: traza ciclo a ciclo y estado final & \repo{pipeline_tb.v} \\
  4. Debug Unit    & Protocolo, carga, modos de ejecución y volcado, desde tramas UART & \repo{tb/unit/debug_*_tb.v}, \repo{top_tb.v} \\
  5. Placa         & Los mismos programas cargados desde la PC, contra el \texttt{.exp} y el modelo & \repo{placa.py}, \repo{comparar.py} \\
  \bottomrule
\end{tabular}
\end{center}

Todos los testbenches son autoverificantes: cuentan los errores, terminan con una línea
\texttt{TEST PASSED} o \texttt{TEST FAILED} y tienen un watchdog para que un procesador que no
llega al HALT no cuelgue la simulación. Los estímulos aleatorios usan semilla fija. Las
simulaciones de RTL corren en xsim (Vivado), porque los testbenches que usan las memorias simulan
el IP de la BRAM.

\subsection{Dos oráculos independientes}

Una prueba vale lo que vale su resultado esperado. Hay dos fuentes, escritas por separado del RTL
y entre sí:

\begin{description}[style=nextline,nosep,leftmargin=1.2em]
  \item[Estado esperado a mano (\repo{sw/<prog>.exp})] Registros, memoria y ciclos al terminar,
        calculados leyendo \repo{isa.md} a partir de un comentario por instrucción en el
        \texttt{.s}. \textbf{Nunca} se copian de la salida del modelo ni del RTL: si salieran de
        alguno de los dos, un error de interpretación de la ISA estaría en ambos lados y la
        prueba pasaría igual.
  \item[Modelo de referencia (\repo{tools/isasim})] Ejecuta cada instrucción y, con las reglas de
        \repo{pipeline.md}, calcula qué hay en cada registro de segmentación en cada ciclo
        (sección~\ref{sec:interfaz}).
\end{description}

El modelo tiene que coincidir con el \texttt{.exp} de todos los programas antes de usarse para
verificar el RTL, y el RTL tiene que coincidir con los dos. Cuando dos fuentes escritas por
separado coinciden, es poco probable que compartan el mismo error.

\subsection{Programas de prueba}

Hay 39 programas en \repo{sw/} (\decision{019}). Cada uno prueba \textbf{un mecanismo}, termina
en \texttt{halt} y tiene un comentario por instrucción con el valor que deja:

\begin{itemize}[nosep]
  \item \textbf{Uno por tipo de instrucción} (10): \repo{alu_r}, \repo{alu_i}, \repo{shifts},
        \repo{lui}, \repo{mem_palabra}, \repo{mem_byte_media}, \repo{branch_cond}, \repo{jump},
        \repo{halt_basico}, \repo{mem_avisos}. Entre todos cubren las 32 instrucciones y el HALT.
  \item \textbf{Uno por caso de riesgo} (12): forwarding desde cada etapa y con prioridad,
        destino \texttt{x0}, bypass del banco, carga-uso en \texttt{rs1} y \texttt{rs2}, load
        seguido de \texttt{sw}, \texttt{beq} y \texttt{jalr}, load sin dependencia, saltos que
        dependen de la instrucción anterior, HALT retenido por un stall, HALT en el camino
        equivocado y un lazo realista (\repo{mixto_lazo}).
  \item \textbf{Variantes sin riesgos} (10 \repo{*_nops}): resuelven los riesgos en software con
        \texttt{nop}, para probar las cinco etapas antes de tener forwarding.
  \item \textbf{Independientes} (7 \repo{indep_*}): parten de un estado inicial y no leen
        registros que escriben, para las integraciones anteriores a WB.
\end{itemize}

Los tests de \repo{sw/} comprueban que toda instrucción aparece en algún programa, que todo caso de
riesgo de \repo{pipeline.md} tiene el suyo y que cada \texttt{.s} tiene su \texttt{.exp}.

\subsection{Banco de pruebas incremental}

\repo{tb/integration/pipeline_tb.v} es un único banco para todas las etapas
(\decision{021}). Carga el programa en la memoria de programa por el puerto~B, como la Debug Unit,
y en \textbf{cada ciclo} compara los cuatro registros de segmentación con la traza del modelo:
\texttt{(dut \^{} ref) \& mask}, donde la máscara ignora los campos de datos de las burbujas. Con
\texttt{+FINAL} compara además el estado final con el \texttt{.exp}; con \texttt{+PAUSAS=<semilla>}
baja \texttt{enable} en ciclos al azar y comprueba que con el core detenido nada cambia.

Las etapas se habilitaron con \emph{defines} acumulativos (\texttt{STAGE\_IF}, \texttt{STAGE\_ID},
\ldots): cada integración comparó solo los latches de las etapas que ya existían, así que cada error
quedó acotado a la última etapa agregada.

\subsection{Regresión y mutantes}

\repo{scripts/regresion.py} y \repo{regresion.tcl} corren \repo{pipeline_tb} con los 39 programas y
después los 27 testbenches unitarios, en una sola simulación de Vivado por lotes (decisiones 028,
032 y 033). La regresión de \texttt{master} pasa \textbf{39 de 39} corridas y \textbf{27 de 27}
testbenches.

Para saber si los programas cubren de verdad cada mecanismo, \repo{scripts/mutantes.py} tiene
\textbf{39 mutantes} del RTL: cada uno rompe una regla de \repo{pipeline.md} (el forwarding sin
prioridad, el stall sin burbuja, el salto sin flush de IF/ID, la parada por HALT mirando solo ID,
\ldots). \textbf{Todos} hacen fallar al menos una corrida (\decision{031}). Los mutantes que detecta
un solo programa son justamente los casos que tienen un programa propio: la prioridad del
forwarding, el destino \texttt{x0}, el HALT en el camino equivocado y el HALT retenido por un
stall. Dos mutantes son equivalentes (ningún programa puede activar esos términos en el core) y los
prueba el testbench de \repo{hazard_unit} con las entradas sueltas.

\subsection{Ciclos por programa}

Con la regresión aprobada, el modelo, el \texttt{.exp} y el RTL dan el mismo número de ciclos, que
cumple
\[
  \text{ciclos} = n + 4 + \text{stalls} + 2 \cdot \text{saltos tomados}
\]
con $n$ las instrucciones que llegan a WB (HALT incluido) y 4 ciclos de llenado. En los 22 programas
del plan:

\begin{center}
\small
\begin{tabular}{@{}l r r r r r@{}}
  \toprule
  & $n$ & \textbf{Stalls} & \textbf{Saltos tomados} & \textbf{Ciclos} & \textbf{CPI} \\
  \midrule
  \repo{loaduse_sw_beq} & 19 & 5 & 2 & 32 & 1,68 \\
  \repo{salto_depende}  & 13 & 0 & 3 & 23 & 1,77 \\
  \repo{mixto_lazo}     & 58 & 5 & 8 & 83 & 1,43 \\
  \midrule
  \textbf{Total de los 22} & 394 & 14 & 28 & 552 & 1,40 \\
  \bottomrule
\end{tabular}
\end{center}

El flush cuesta más que el stall: 56 ciclos (10\,\%) contra 14 (3\,\%). Sin los 4 de llenado, que
pesan mucho en programas cortos, el CPI es $(552 - 88)/394 = 1{,}18$. Las variantes con
\texttt{nop} tardan entre 2 y 30 ciclos más: \repo{mixto_lazo_nops} tarda 113 contra 83, un 36\,\%
más. Es lo que ahorran en hardware el forwarding y el stall. La tabla completa está en
\repo{verificacion.md} §8.8.

\subsection{Prueba en la placa}

Con el bitstream de \qty{92}{\mega\hertz}:

\begin{itemize}[nosep]
  \item \repo{placa.py} carga los 22 programas del plan uno tras otro (cada uno con un
        \texttt{LOAD}, sin resetear la placa), los corre con \texttt{RUN} y con \texttt{STEP}, y
        compara el estado final con el \texttt{.exp} y cada paso con la traza del modelo:
        \textbf{22 de 22}. Con las variantes \repo{_nops}, \textbf{32 de 32}.
  \item \repo{comparar.py --cada-ciclo} compara un programa con el modelo en todos los ciclos. Los
        dos programas de la demostración (sección~\ref{sec:demo}) coinciden en los 15 y los 44
        ciclos.
\end{itemize}

\captura[0.32\textheight]{comparar_alu_r_y_mixto_lazo.png}{Placa contra el modelo}{Salida de
\repo{comparar.py} con \repo{alu_r} y \repo{mixto_lazo}: el volcado de la placa coincide con la
foto del modelo de referencia.}{fig:comparar}
